\documentclass[12pt,fleqn]{article}
\usepackage[utf8]{inputenc}
\usepackage{paralist} 
\usepackage{amssymb}  
\usepackage{amsthm}   
\usepackage{eurosym} 
\usepackage{multicol} 
\usepackage[left=1.5cm,right=1.5cm,top=1.5cm,bottom=0.5cm]{geometry} 
\usepackage{fancyheadings} 
\pagestyle{fancy} 
\headheight1.6cm 
\lhead{\today} 
\chead{MA Stz v3} 
\rhead{\includegraphics[scale=0.5]{logo.png}} 
\lfoot{} 
    \cfoot{} 
    \rfoot{} 
\usepackage{amsmath}  
\usepackage{cancel} 
\usepackage{pgf,tikz} 
\usetikzlibrary{arrows} 
\newtheoremstyle{aufg} 
{16pt}  % Platz zwischen Kopf und voherigem Text 
{16pt}  % und nachfolgendem Text 
{}     % Schriftart des Koerpers 
{}     % mit \parindent Einzug 
{\bf}  % Schriftart des Kopfes 
{:}     % Nach Bedarf z.B. Doppelpunkt nach dem Kopf 
{0.5em} % Platz zwischen Kopf und Koerper 
{}     % Kopfname 

\theoremstyle{aufg} 
\newtheorem{aufgabe}{Aufgabe} 



\newtheoremstyle{bsp} 
{16pt}  % Platz zwischen Kopf und voherigem Text 
{16pt}  % und nachfolgendem Text 
{}     % Schriftart des Koerpers 
{}     % mit \parindent Einzug 
{\em}  % Schriftart des Kopfes 
{:}     % Nach Bedarf z.B. Doppelpunkt nach dem Kopf 
{0.5em} % Platz zwischen Kopf und Koerper 
{}     % Kopfname 

\theoremstyle{bsp} 
\newtheorem{beispiel}{Beispiel} 



\begin{document} 
    \begin{flushleft}
\renewcommand{\arraystretch}{2.15} 
\begin{tabular}{|p{10cm}|p{2cm}|p{2cm}|p{2cm}|} 
\hline 
\hspace{2cm} Punkte von \qquad20\qquad Punkten erreicht & \hspace{1.2cm} NP & G & E \\ 
\hline 
\end{tabular} \\[1em]    
\begin{center}{\Large Schriftliche Arbeit}\end{center} 
{\bf Zeit: }20 Minuten\\ 
\begin{aufgabe}\framebox{\qquad/10} ~ \\ 
\begin{enumerate}[a)] 
\item 
$2x+2x$=
\item 
$3.39a-3.21a$=
\item 
$4v+2v-3v$=
\item 
$5x+2y-2y-2x-7x$=
\item 
$6.25c-4.6b+5.97c+2.63c+6.47b-3.12a-3.98a-3.7b+2.61a$=
\end{enumerate} 
\end{aufgabe} 
\begin{aufgabe}\framebox{\qquad/10} ~ \\ 
\begin{enumerate}[a)] 
\item 
$6\cdot(4x-4)$=
\item 
$(2+5x)\cdot2$=
\item 
$-3\cdot(6a-2)$=
\item 
$(5a-4b)\cdot5c$=
\item 
$(2a-3b)\cdot 3-2\cdot (-1,5a+4b)=$
\end{enumerate} 
\end{aufgabe} 
\end{flushleft} 
    \end{document}